\documentclass[11pt]{article}
\usepackage[a4paper,margin=1in]{geometry}
\usepackage{amsmath,amssymb}
\usepackage{hyperref}
\title{Harmonic ODF Reconstruction in PyTex}
\author{PyTex Contributors}
\date{\today}

\begin{document}
\maketitle

\section{Scope}

This note documents the band-limited harmonic ODF reconstruction implemented in PyTex for
pole-figure inversion. It is the canonical mathematical companion to the public
\texttt{HarmonicODF} API.

\section{Reference Posture}

Randle and Engler describe the classical series-expansion route by expressing the
orientation distribution function in symmetrized generalized spherical harmonics and the
pole figures in matching harmonic coefficients linked by a linear relation. They also
note the practical effects of truncation and the missing odd coefficients in diffraction
pole figures (\textit{Introduction to Texture Analysis}, pp.\ 105--107).

PyTex follows that scientific posture while keeping the implementation aligned with its
existing explicit frame, symmetry, phase, and pole-figure data model.

\section{Band-Limited Basis}

PyTex uses the Bunge Euler ordering
\[
  (\phi_1, \Phi, \phi_2)
\]
and the standard Wigner-\(D\) building block
\[
  D^l_{mn}(\phi_1,\Phi,\phi_2)
  =
  e^{- i m \phi_1}
  d^l_{mn}(\Phi)
  e^{- i n \phi_2}.
\]

The implemented real basis is constructed from
\[
  \Re D^l_{mn}
  =
  d^l_{mn}(\Phi)\cos(m\phi_1 + n\phi_2),
\]
\[
  \Im D^l_{mn}
  =
  d^l_{mn}(\Phi)\sin(m\phi_1 + n\phi_2),
\]
truncated at a user-selected degree bandlimit \(l_{\max}\).

\section{Symmetry Handling}

The ODF must satisfy
\[
  f(g h) = f(g)
  \qquad \forall h \in C
\]
for crystal symmetry and
\[
  f(s g) = f(g)
  \qquad \forall s \in S
\]
for specimen symmetry.

PyTex enforces these invariance relations by explicit group averaging of each basis
function:
\[
  \bar{\Psi}(g)
  =
  \frac{1}{|S||C|}
  \sum_{s \in S}
  \sum_{h \in C}
  \Psi(s g h).
\]

This left/right split is not optional bookkeeping. It is the mechanism that keeps the
harmonic surface consistent with the canonical PyTex frame and symmetry model.

\section{Quadrature And Orthonormalization}

PyTex evaluates the projected basis on a full Bunge grid over \(SO(3)\) using midpoint
quadrature and Haar weights proportional to
\[
  \sin\Phi.
\]

Let the quadrature orientations be \(g_q\) with normalized weights \(w_q\). The raw
projected basis matrix is
\[
  B_{q\alpha} = \bar{\Psi}_{\alpha}(g_q).
\]

PyTex forms the weighted Gram matrix
\[
  G = B^{\mathsf T} W B
\]
with \(W = \mathrm{diag}(w_q)\), then retains the numerically non-singular eigenspace and
constructs an orthonormal basis on the quadrature grid.

\section{Forward Operator}

PyTex does not introduce a second incompatible pole-figure object for harmonic inversion.
Instead it uses the current measured-direction \texttt{PoleFigure} semantics and builds a
regularized pole-density operator against the harmonic ODF sampled on the quadrature grid:
\[
  I_m
  \approx
  \sum_q
  w_q f(g_q) A_{mq}.
\]

For a pole family \(H\) and a measurement direction \(y_m\),
\[
  A_{mq}
  =
  \frac{1}{|H|}
  \sum_{u \in H}
  K\!\left(\angle(y_m, g_q u)\right),
\]
where \(K\) is the configured pole-density kernel.

\section{Inverse Problem}

Writing the band-limited ODF as
\[
  f(g) \approx \sum_{\alpha} c_{\alpha}\Phi_{\alpha}(g),
\]
PyTex solves the regularized least-squares problem
\[
  \min_{\mathbf{c}}
  \frac{1}{2}
  \lVert A\mathbf{c} - \mathbf{b} \rVert_2^2
  +
  \frac{\lambda}{2}
  \lVert \mathbf{c} \rVert_2^2.
\]

The implementation uses dense linear least squares on the augmented system. This is
appropriate for the present band-limited and explicitly auditable scientific scope.

\section{Even-Degree Handling}

When all input pole figures are antipodal, PyTex defaults to retaining even degrees only.
This follows the classical diffraction-texture limitation that odd coefficients are not
uniquely observable from antipodal pole figures.

\section{Interpretation Notes}

\begin{itemize}
  \item Harmonic reconstruction is now implemented, but it remains a band-limited model.
  \item Symmetry handling is explicit through the repository point-group operators, not a
        separate harmonic-specific symmetry abstraction.
  \item Negative ODF lobes can still appear because truncated harmonic expansions are not
        positivity preserving.
  \item Degree selection, quadrature density, and kernel halfwidth are scientific inputs,
        not hidden tuning constants.
\end{itemize}

\section{References}

\begin{thebibliography}{9}
\bibitem{randle}
V. Randle and O. Engler,
\textit{Introduction to Texture Analysis: Macrotexture, Microtexture and Orientation Mapping},
CRC Press, 2000.

\bibitem{bunge}
H.-J. Bunge,
\textit{Texture Analysis in Materials Science: Mathematical Methods},
Butterworths, 1969.

\end{thebibliography}

\end{document}
